% =====================================================================
%  sections/sec1.tex —— 1 引言（原文 pp.1–3）
% =====================================================================
\section{引言}
\label{sec:intro}

对共形场论（CFT）的研究，在理解量子场论动力学上起着重要作用。最著名的四维 CFT 是 $\sN{=}4$ 超对称规范理论：它的共形流形由严格边缘耦合 $\tau$（即规范耦合常数）标记。带显式拉氏量的弱耦合描述，只能写在耦合常数空间的尖点处。不同的弱耦合描述之间由 S 对偶联系；从 magical 六维 $(2,0)$ 理论在环面上的紧致化出发，最容易看清这一点 \cite{1}。

\zhnote{[译者按（无外部依据）：magical 是这一领域对六维 $(2,0)$ 理论的一种戏称，作者借此标记其性质之奇特，它不是技术术语，故保留原文不译。]}

近期，人们用带正则穿刺点的黎曼面紧致化六维 $(2,0)$ 理论，找到了很大一类四维 $\sN{=}2$ 超共形场论（SCFT）\cite{2,3}。与此类似，规范耦合常数空间等同于带穿刺黎曼面的复结构模空间，S 对偶群等同于模群。然而，尖点处的弱耦合对偶框架通常写不出拉氏量描述。原因在于：物质部分往往是孤立强耦合 SCFT，而它们在理解这些理论的 S 对偶时起关键作用。例如，带六个基本表示的 $\SU(3)$ 的 S 对偶理论 \cite{4} 牵涉强耦合的 $E_6$ 理论 \cite{5}，后者可用六维 $A_2$ 理论在带三个满穿刺点的球面上紧致化来构造。

上面这些理论的耦合常数无量纲，库仑分支算子有整数标度维。还有另一类 $\sN{=}2$ SCFT，称为 Argyres--Douglas（AD）理论 \cite{6}：其库仑分支算子通常带分数标度维，耦合常数有量纲。这类理论最初是作为纯 $\SU(3)$ 规范理论库仑分支上某个特殊点的红外理论被发现的。在该特殊点上，彼此非局域的双荷子变成无质量，于是不存在使这些双荷子只带电荷的对偶框架。这个理论必定是相互作用的 SCFT \cite{6,7}，而且不可能有拉氏量描述。它的谱中通常出现相关算子，可以用这些算子把理论形变到新的不动点。

自然要问：能否用六维 $(2,0)$ 理论来工程化 AD 理论，并把这些物理性质全部编码到黎曼面上的几何对象里？答案是肯定的，但必须引入新类型的穿刺点：非正则奇点（高阶极点）\cite{3}。此前的分析以 $A_1$ $(2,0)$ 理论为基础；本文将针对更高秩的情形做完整的分析。由于极点的阶可以趋于无穷，找到新理论的可能性立刻大大增加。而一旦找到六维构造，分类与研究这些理论就简单得惊人。下面我们先勾勒构造 AD 理论的主要思路，并概述主要结果。

\zhnote{[来源：复常微分方程教材中 regular singular point 与 irregular singular point 的标准定义；I：本库 SY03-04 冻结表作「非正则穿刺点」。]「非正则」是相对于「正则奇点」而言的：正则奇点指极点阶不超过一的奇点，此处泛指极点阶大于一的奇点。它说的是奇点邻域内解的行为，而不是奇点的形状，故不作「不规则奇点」。}

定义在黎曼面上的 Hitchin 方程 \cite{8,9} 在这些构造中起核心作用。先回顾正则穿刺点的情形：四维 UV 理论由带穿刺点的黎曼面指定，而穿刺点被解释为 Hitchin 方程中场的边界条件——这些穿刺点具有正则奇点（一阶极点）。规范耦合常数等同于带穿刺黎曼面的复结构模，质量参数编码为一阶极点的系数。四维理论的红外行为由满足上述边界条件的解模空间决定。特别地，Seiberg--Witten 曲线等同于 Hitchin 可积系统的谱曲线。

只用正则穿刺点，无法引入有量纲的耦合，因此需要 Hitchin 方程解的非正则奇点。引入非正则奇点恰好提供我们想要的性质：第一，四维理论的库仑分支算子可以有分数标度维，其数值由谱曲线算出；第二，高阶极点中的参数就是有量纲的耦合常数。一阶极点的系数仍然等同于质量参数，于是 UV 理论的全部形变参数都与几何参数对上了。模空间的性质与正则奇点的情形类似，即它同样是超凯勒流形 \cite{10}。我们也可以把模空间的谱曲线等同于 Seiberg--Witten 曲线，理论的红外行为同样得到求解。

正则奇点由杨图分类；可以在任意亏格的黎曼面上放任意多个穿刺点来定义 SCFT。非正则奇点的情形则完全不同。首先，黎曼面的坐标 $z$ 在四维 $U(1)_R$ 对称性下按非平凡方式变换，为保持 $U(1)_R$ 对称性，只能使用黎曼球。其次，能放在同一黎曼面上的非正则奇点与正则奇点的组合，还受到很严格的限制。最后，非正则奇点的分类结果丰富得多，而且并非每个非正则奇点都定义 AD 理论。要构造 AD 理论，有以下约束：
\begin{enumerate}[label=(\arabic*)]
  \item 只能用黎曼球。
  \item 只有两种奇点组合：(a) 只有一个非正则奇点；(b) 一个非正则奇点加一个正则奇点。
  \item 非正则奇点只有三类。
\end{enumerate}

只要指定黎曼球上的奇点结构，就能写出谱曲线，从而定出理论中各种算子的标度维数。部分 AD 理论的三维镜像理论也可以由非正则奇点的信息确定。用三维镜像理论可以轻松算出这些理论的中心荷 $a$ 与 $c$。我们构造出的理论涵盖了文献中几乎所有的 AD 理论，此外还有大量新例子。构造出的 SCFT 并非全部互不相同；我们识别出一些有用的同构，它们对计算这些理论的中心荷等量有用。

即使在上述约束下，允许高阶极点也大大增加了构造 SCFT 的可选择余地。我们的结果极大扩展了 $\sN{=}2$ SCFT 的版图，并表明这些 AD 理论比人们以为的要普遍得多。事实上，定义在球面上的 class S 理论（即用正则奇点的理论）也可以用非正则奇点来实现。

本文还研究了这些理论的一些其他性质。第一，我们用奇点的碰撞来识别 $\SU(N)$ QCD 的 AD 点。第二，如果 AD 理论带一个正则奇点（它通常带非阿贝尔味对称性），就可以用它构造新的渐近自由（AF）理论。几何上，这类 AF 理论由在黎曼面上放任意多个非正则奇点来工程化。

本文的组织如下：第 \ref{sec:general} 节简要回顾 AD 理论；第 \ref{sec:hitchin} 节给出 Hitchin 方程非正则奇点的分类；第 \ref{sec:a1} 节讨论由六维 $A_1$ 理论得到的 AD 理论，这些理论并不新，但我们所研究的各种性质似乎是新的；第 \ref{sec:a2} 节分类并研究由 $A_2$ 理论得到的 AD 理论；第 \ref{sec:ak} 节考虑由紧致化六维 $A_{k-1}$ 理论得到的一般 AD 理论；第 \ref{sec:known} 节讨论已知 AD 理论的六维构造；最后在第 \ref{sec:concl} 节给出一个简短的讨论，指出可能的研究方向。
